\wordchapter{1}{Introduction}

Weather shocks can cause losses across many farms at once, limiting the scope for private insurers to pool risk. Insurers also have incomplete information about farm risks and farmers' efforts to prevent losses (Miranda \& Glauber, 1997; Mahul \& Stutley, 2010). South Korea supports crop disaster insurance through premium subsidies, administrative cost subsidies, and state reinsurance. The 2026 implementation guidelines list 78 crops and a budget of KRW 556.6 billion for national premium and loading support (Ministry of Agriculture, Food and Rural Affairs [MAFRA], 2026a).

These market constraints help explain government intervention, but the design of support also depends on its objectives and incentives. Korean studies examine production and price effects (Han, 2014), input use and farm income (Nam \& Ahn, 2022), and output and technical efficiency among fruit farms (Lee \& Kim, 2026). This evidence describes agricultural outcomes. Evaluating the support instruments also requires examining how risks are shared, who receives assistance, and what public expenditure achieves.

In this thesis, I examine the roles of the three support instruments and the production patterns that accompany insurance availability. I first assess financing arrangements and incentives using program rules and financing data. I then compare six annual crops with crops not yet exposed to insurance and compare rice planted area in early and later pilot municipalities. For the national analysis, I separate changes in area from changes in yield and examine alternative reference periods, comparison crops, and aggregation weights.

The institutional assessment identifies different evaluation priorities: access and distribution for premium support, delivery costs and service quality for administrative support, and risk allocation and pricing for state reinsurance. The national production contrast is −0.002 log points, but estimates vary across specifications and pre-pilot growth differs across groups. The rice area estimate is −0.021 log points and changes with sample composition. Both estimates have intervals that include zero. The rice comparison holds crop identity fixed and covers the first two years of availability. These comparisons concern insurance availability as a policy package; the institutional analysis examines the roles of the individual instruments.

Chapters 2 and 3 describe the program and assess the three forms of government support. Chapters 4 and 5 present the data, empirical strategy, and results. Chapters 6 and 7 discuss the policy implications and conclude.
